%-------------------------
% Resume in Latex
% Author : Jake Gutierrez
% Based off of: https://github.com/sb2nov/resume
% License : MIT
%------------------------

\documentclass[letterpaper,11pt]{article}

\usepackage{latexsym}
\usepackage[empty]{fullpage}
\usepackage{titlesec}
\usepackage{marvosym}
\usepackage[usenames,dvipsnames]{color}
\usepackage{verbatim}
\usepackage{enumitem}
\usepackage[hidelinks]{hyperref}
\usepackage{fancyhdr}
\usepackage[english]{babel}
\usepackage{tabularx}
\usepackage{fontawesome5}
\usepackage{graphicx}
\usepackage[normalem]{ulem}
\input{glyphtounicode}
\usepackage{enumitem}
\usepackage{xcolor}
\usepackage{ragged2e}


%----------FONT OPTIONS----------
\pagestyle{fancy}
\fancyhf{} 
\fancyfoot{}
\renewcommand{\headrulewidth}{0pt}
\renewcommand{\footrulewidth}{0pt}

% Adjust margins
\addtolength{\oddsidemargin}{-0.5in}
\addtolength{\evensidemargin}{-0.5in}
\addtolength{\textwidth}{1in}
\addtolength{\topmargin}{-0.5in}
\addtolength{\textheight}{1.0in}

\urlstyle{same}

\raggedbottom
\raggedright
\setlength{\tabcolsep}{0in}

% Sections formatting
\titleformat{\section}{
  \vspace{-5pt}\scshape\raggedright\large
}{}{0em}{}[\color{black}\titlerule \vspace{-7pt}]

% Ensure that generate pdf is machine readable/ATS parsable
\pdfgentounicode=1

%-------------------------
% Custom commands
\newcommand{\resumeItem}[1]{
  \item\small{
    {#1 \vspace{-2pt}}
  }
}

\newcommand{\resumeSubheading}[4]{
  \vspace{-2pt}\item
    \begin{tabular*}{0.97\textwidth}[t]{l@{\extracolsep{\fill}}r}
      \textbf{#1} & #2 \\
      \textit{\small#3} & \textit{\small #4} \\
    \end{tabular*}\vspace{-7pt}
}

\newcommand{\resumeSubSubheading}[2]{
    \item
    \begin{tabular*}{0.97\textwidth}{l@{\extracolsep{\fill}}r}
      \textit{\small#1} & \textit{\small #2} \\
    \end{tabular*}\vspace{-7pt}
}

\newcommand{\resumeProjectHeading}[2]{
    \item
    \begin{tabular*}{0.97\textwidth}{l@{\extracolsep{\fill}}r}
      \small#1 & #2 \\
    \end{tabular*}\vspace{-6pt}
}

\newcommand{\resumeSubItem}[1]{\resumeItem{#1}\vspace{-4pt}}

\renewcommand\labelitemii{$\vcenter{\hbox{\tiny$\bullet$}}$}

\newcommand{\resumeSubHeadingListStart}{\begin{itemize}[leftmargin=0.15in, label={}]}
\newcommand{\resumeSubHeadingListEnd}{\end{itemize}\vspace{-10pt}} 
\newcommand{\resumeItemListStart}{\begin{itemize}}
\newcommand{\resumeItemListEnd}{\end{itemize}\vspace{-7pt}}

%------------------------------------------
%%%%%%  RESUME STARTS HERE  %%%%%%%%%%%%%%%%%%%%%%%%%%%%


\begin{document}



%----------HEADING----------
\begin{center}
    \textbf{\Huge \scshape Satish Shrishail Mirji} \\ \vspace{5pt}
    \normalsize
    \faMobile*\ \underline{\href{tel:+91-9380240695}{\textcolor{blue}{+91-9380240695}}} \quad
    \faEnvelope\ \underline{\href{mailto:satish15.dev@gmail.com}{\textcolor{blue}{satish15.dev@gmail.com}}} \quad
    \faLinkedin\ \underline{\href{https://www.linkedin.com/in/satish-m-a99b66334/}{\textcolor{blue}{LinkedIn}}} \quad
    \faGithub\ \underline{\href{https://github.com/techwithsatish}{\textcolor{blue}{GitHub}}} \quad
    \faBook\ \underline{\href{https://satish15.hashnode.dev/}{\textcolor{blue}{Hashnode}}}
\end{center}

\vspace{-5pt}




%-----------TECHNICAL SKILLS-----------

\section{Technical Skills}

\begin{tabular}{|p{6cm}|p{12cm}|}
\hline
\textbf{Cloud Technologies} & AWS \\ \hline
\textbf{Build Tools} & Gradle, Maven, NPM \\ \hline
\textbf{Container Technologies} & Docker, Kubernetes (Amazon EKS), Docker Compose \\ \hline
\textbf{Scripting} & Python, Shell Scripting (Bash) \\ \hline
\textbf{Version Control} & Git, GitHub \\ \hline
\textbf{Configuration Management} & Ansible \\ \hline
\textbf{Infrastructure as Code} & Terraform \\ \hline
\textbf{CI/CD} & GitHub Actions, Jenkins, Argo CD \\ \hline
\textbf{Monitoring \& Observability} & Prometheus, Grafana, \\ \hline
\textbf{Operating Systems} & Linux (Ubuntu) \\ \hline
\textbf{Networking} & VPC, IAM, Route 53, Application Load Balancer (ALB) \\ \hline
\end{tabular}


%-----------EDUACTION-----------
\section{Education}
  \resumeSubHeadingListStart
    \resumeSubheading
      {Sai Vidya Institute Of Technology Bengaluru}{2022 - 2026}
      {B.E in CSE (Data Science) - CGPA: 8.3}{Karnataka}
      
    \vspace{2pt}
      
    \resumeSubheading
      {Sainik School Kodagu}{2021 }
      {CBSE Class 12th - Percentage: 80\%}{Karnataka}

    \vspace{2pt}

    \resumeSubheading
      {Sainik School Kodagu}{2019 }
      {CBSE Class 10th - Percentage: 84\%}{Karnataka}
  \resumeSubHeadingListEnd
\vspace{1mm}

%-----------WORK EXPERIENCE-----------


\section{Work Experience}
\justifying

\begin{itemize}[leftmargin=0.15in, label={}]
\small{\item{
\textbf{Full Stack (MERN) Developer Intern} \hfill Dec 2025 -- May 2026 \\
\textit{SuprMentr, Bengaluru, Karnataka, India} \\
\vspace{-20pt}

\begin{itemize}
    \justifying

    \item Worked as a \textbf{Full Stack (MERN) Developer Intern} and developed the
    \textbf{\href{https://github.com/techwithsatish/ecommerce-website}{\textcolor{blue}{\underline{Ecommerce Website (MERN Stack)}}}}
    using \textbf{MongoDB, Express.js, React.js, and Node.js} as Personal Project.

    \item Built a dynamic and mobile-responsive user interface using \textbf{React.js} and \textbf{Tailwind CSS}, implementing \textbf{product filtering}, \textbf{product variations}, and \textbf{real-time cart state management}.

    \item Designed and developed \textbf{RESTful APIs} using \textbf{Node.js} and \textbf{Express.js} for \textbf{user authentication}, \textbf{product management}, and \textbf{order processing}.

    \item Designed the database schema using \textbf{MongoDB} and \textbf{Mongoose} to efficiently manage \textbf{user profiles}, \textbf{product catalogs}, and \textbf{order history}.

    \item Integrated the \textbf{Stripe Payment Gateway} to support secure \textbf{online payments} and \textbf{Cash on Delivery (COD)}.

    \item Developed a dedicated \textbf{Admin Dashboard} for \textbf{inventory management}, including \textbf{product upload} (with image handling), \textbf{product deletion}, and \textbf{order tracking}.

    \item Deployed the application on \textbf{Vercel}, configuring \textbf{frontend-backend integration}, environment variables, and production routing.

\end{itemize}

}}
\end{itemize}


\vspace{-13pt}




%-----------PROJECTS-----------

\section{Projects}
\justifying
\resumeSubHeadingListStart

\resumeProjectHeading
{\textbf{\underline{E-COMMERCE DEVOPS Implementation}}}{}
\resumeItemListStart

\resumeItem{Implemented DevOps practices such as \textbf{Docker, Kubernetes, Infrastructure as Code (Terraform), and CI/CD} for a real-world multi-microservice e-commerce application.}

\resumeItem{Deployed and managed a \textbf{multi-microservice architecture} on \textbf{Amazon EKS}, gaining hands-on experience with container orchestration and Kubernetes workloads.}

\resumeItem{Provisioned AWS infrastructure including \textbf{VPC, Subnets, Security Groups, IAM Roles, and Amazon EKS} using reusable \textbf{Terraform modules}.}

\resumeItem{Implemented an automated \textbf{CI/CD pipeline} using \textbf{GitHub Actions} and \textbf{Argo CD} to achieve continuous integration and GitOps-based deployments.}

\resumeItem{Configured AWS networking services including \textbf{VPC, Application Load Balancer (ALB), Route 53, and IAM} to securely expose the application.}

\resumeItem{Worked with \textbf{Docker Compose}, Kubernetes manifests, and cloud-native deployment practices to automate application deployment and management.}

\resumeItemListEnd

\resumeProjectHeading
{\textbf{\underline{AWS Cloud Cost Optimization using AWS Lambda}}}{}

\resumeItemListStart

\resumeItem{Developed a \textbf{serverless AWS cost optimization solution} using \textbf{AWS Lambda}, \textbf{Python}, and \textbf{Boto3} to automatically identify and remove stale \textbf{Amazon EBS snapshots}, reducing unnecessary cloud storage costs.}

\resumeItem{Integrated the \textbf{AWS Boto3 SDK} with \textbf{EC2}, \textbf{EBS}, and AWS APIs to retrieve instances, volumes, and snapshot metadata, implementing logic to detect orphaned resources while preserving active infrastructure.}

\resumeItem{Configured secure \textbf{IAM Roles and Policies} following the \textbf{Principle of Least Privilege (PoLP)} to enable automated snapshot management and infrastructure maintenance.}

\resumeItem{Automated periodic execution using \textbf{Amazon CloudWatch (EventBridge)} to eliminate manual cloud resource cleanup, improving cloud governance and optimizing AWS infrastructure costs.}

\resumeItemListEnd

\resumeProjectHeading
{\textbf{\underline{GitHub–Jira Automation using Python and Flask}}}{}

\resumeItemListStart

\resumeItem{Developed an automation solution using \textbf{Python}, \textbf{Flask}, and \textbf{GitHub Webhooks} to automatically create \textbf{Jira} issues from validated GitHub issue comments, eliminating manual ticket creation.}

\resumeItem{Integrated \textbf{GitHub Webhooks} with a \textbf{Flask REST API} hosted on an \textbf{AWS EC2} instance to receive real-time JSON payloads and trigger workflow automation.}

\resumeItem{Implemented \textbf{Jira REST API} integration using the \textbf{Requests} library with secure API token authentication to create project-specific issues, populate metadata, and automate issue tracking.}

\resumeItem{Built conditional request processing to validate GitHub issue comments before creating Jira tickets, improving workflow efficiency and reducing manual operational effort for development teams.}

\resumeItemListEnd
\resumeProjectHeading
{\textbf{\underline{GitHub Repository Access Automation using Shell Scripting}}}{}

\resumeItemListStart

\resumeItem{Developed a \textbf{Bash Shell Script} to automate retrieval of GitHub repository collaborators using the \textbf{GitHub REST API}, eliminating repetitive manual access verification tasks.}

\resumeItem{Integrated the script with the \textbf{GitHub REST API} using \textbf{cURL} and Personal Access Token (PAT) authentication to securely fetch repository user and permission details.}

\resumeItem{Parsed JSON API responses using \textbf{jq} to filter collaborator permissions and identify users with repository access, improving administrative visibility and access management.}

\resumeItem{Implemented reusable shell functions, environment variable configuration, command-line arguments, and validation logic to create a scalable Linux automation solution for repository administration.}

\resumeItemListEnd

\resumeSubHeadingListEnd



%-----------CERTIFICATIONS-----------
\section{Certifications}
\begin{itemize}[leftmargin=0.15in, label={}]
    \small{\item{

\textbf{\href{https://www.udemy.com/certificate/UC-906733f7-7512-4649-972f-e33253762d3e/}{\color{blue}Full Stack Web Development Certification – Udemy}} \\ \vspace{-15pt}

\textbf{DevOps Zero to Hero – Abhishek Veeramalla (YouTube)} \\ \vspace{-15pt}

    }}
\end{itemize}
\vspace{-13pt}

\end{document}
